% ===== 原文 pp.37–43（PDF pp.38–44）=====
% tr-05 片段 a：原文 p.37（PDF p.38）
% 内容：§4 结尾 (4.55)–(4.56) → \section{镜像对称} → 5.1 荷电场的 T-对偶 → (5.1)–(5.4)
% [原文 p.37]

这意味着：若 $\phi_i=0$ 所对应的那张边界超平面存在，则 $v^i$ 是多面体 $P_\zeta\subset \ii\ft_F^*$ 在该边界超平面处的内向法向。对子集 $I\subset\{1,\dots,N\}$，令 $\sigma_I$ 为 $\{v^i\}_{i\in I}$ 张成的锥，并约定 $\sigma_\emptyset:=\{0\}$。于是 $X_\zeta$ 就是相应于扇

\begin{equation}
    \Delta_\zeta=\Bigl\{\;\sigma_I\;\Bigm|\; I\subset\{1,\dots,N\},\ I\notin\cI_\zeta\;\Bigr\}.
    \tag{4.55}
\end{equation}
的环面簇。特别地，$\Delta_\zeta$ 中的一维锥是

\begin{equation}
    \Delta_\zeta^{(1)}=\Bigl\{\;\R_{\ge0}v^i\;\Bigm|\; \exists\ \text{对应于 }\phi_i=0\text{ 的 }P_\zeta\text{ 边界}\;\Bigr\}.
    \tag{4.56}
\end{equation}

% [Tutor Note: p.37 式 (4.55)(4.56) 中的向量写作 $v^i$（上标），文字层会塌成 $vi$，抄写时极易误作下标 $v_i$；判据是 (4.56) 里同一个符号出现在 $\R_{\ge0}v^i$ 中作为锥的生成元，而上标形式与 (4.55) 的 $\{v^i\}_{i\in I}$ 一致。另：$\cI_\zeta$ 是花体 $\mathcal{I}$（文字层抽成普通 $I$），(4.55) 里的条件 $I\notin\cI_\zeta$ 与正文中“若该边界超平面存在”呼应。]

\section{镜像对称}

\subsection{荷电场的 T-对偶}

基本想法是对每个物质多重态中标量分量的相位施加 T-对偶。为此，我们先考虑这样一个理论：其中只有一个标量 $\phi$，它在 $U(1)$ 规范对称性下带荷 $1$。用极坐标变量 $(\rho,\varphi)$（由 $\phi=\rho\,\ee^{\ii\varphi}$ 定义）表出时，标量动能项 $|D_0\phi|^2-|D_1\phi|^2$ 等于 $(\partial_0\rho)^2-(\partial_1\rho)^2$ 再加上

\begin{equation}
    \cL_\varphi=\rho^2(\partial_0\varphi+v_0)^2-\rho^2(\partial_1\varphi+v_1)^2.
    \tag{5.1}
\end{equation}

% [Tutor Note: p.37 原文 “let us first consider the theory is a single scalar $\phi$” 语法不通（应为 “the theory with a single scalar $\phi$” 或 “the theory is a single scalar $\phi$”），译文按“只有一个标量 $\phi$ 的理论”处理；这不影响 (5.1) 之后的任何推导。另注意本页起 $\phi$（标量场）与 $\varphi$（其相位）是两个字形的不同符号，不能混用。]

我们现在用路径积分的论证来对 $\varphi$ 施加 T-对偶。考虑由 $\varphi:\Sigma\to\R/2\pi\Z$ 与一个一形式 $J\in\Omega^1(\Sigma,\R)$ 组成的系统，其拉氏量为

\begin{equation}
    \cL'=\frac{1}{4\rho^2}\bigl(J_0^2-J_1^2\bigr)+J_1(\partial_0\varphi+v_0)-J_0(\partial_1\varphi+v_1).
    \tag{5.2}
\end{equation}

若先积掉 $J$-场（其对 $J$ 的极值点为 $J_0=2\rho^2(\partial_1\varphi+v_1)$ 与 $J_1=2\rho^2(\partial_0\varphi+v_0)$），我们便回到 (5.1) 中的拉氏量 $\cL_\varphi$。另一方面，若先积掉 $\varphi$，则得到约束 $J_\mu=\partial_\mu\tphi$，其中 $\tphi:\Sigma\to\R/2\pi\Z$ 是某个映射。把它代入拉氏量的其余部分，得到

\begin{equation}
\begin{aligned}
    \cL_{\tphi}&=\frac{1}{4\rho^2}\bigl((\partial_0\tphi)^2-(\partial_1\tphi)^2\bigr)+\partial_1\tphi\,v_0-\partial_0\tphi\,v_1\\
    &\simeq\frac{1}{4\rho^2}\bigl((\partial_0\tphi)^2-(\partial_1\tphi)^2\bigr)+\tphi\,v_{01},
\end{aligned}
    \tag{5.3}
\end{equation}

其中 $\simeq$ 表示在相差一个全导数的意义下相等。我们看到，对偶变量 $\tphi$ 作为一个动力学的 theta 角进入作用量。注意 $(\rho,\tphi)$ 的度规由

\begin{equation}
    \dd\tilde{s}^2=\dd\rho^2+\frac{1}{4\rho^2}\dd\tphi^2=\frac{|\dd(\rho^2)+\ii\dd\tphi|^2}{4\rho^2}=\frac{|\dd y|^2}{2(y+\overline{y})},
    \tag{5.4}
\end{equation}
给出。

% [Tutor Note: p.37 式 (5.2)–(5.4) 可逐式自检：对 $J_0$ 变分得 $J_0=2\rho^2(\partial_1\varphi+v_1)$、对 $J_1$ 变分得 $J_1=2\rho^2(\partial_0\varphi+v_0)$，代回完整的 $L'$：二次项给 $\rho^2[(\partial_1\varphi+v_1)^2-(\partial_0\varphi+v_0)^2]$，交叉项 $J_1(\partial_0\varphi+v_0)-J_0(\partial_1\varphi+v_1)$ 给 $2\rho^2[(\partial_0\varphi+v_0)^2-(\partial_1\varphi+v_1)^2]$，两者合计才是 $\cL_\varphi$；对 $\varphi$ 积分给出 $J_\mu=\partial_\mu\tphi$，再对 $J_1v_0-J_0v_1$ 分部积分即得 $v_{01}:=\partial_0v_1-\partial_1v_0$。度规一项的判据：$y=\rho^2+\ii\tphi$ 时 $\dd y=\dd(\rho^2)+\ii\dd\tphi$、$y+\overline{y}=2\rho^2$，故 $|\dd y|^2/(2(y+\overline{y}))=\bigl((\dd\rho^2)^2+\dd\tphi^2\bigr)/(4\rho^2)=\dd\rho^2+\frac{1}{4\rho^2}\dd\tphi^2$，与 (5.4) 前两项逐项相等。]

% [Tutor Note: p.37 式 (5.4) 左边的度规是 $\dd\tilde s^2$（$s$ 上有波浪号，文字层抽成 “des2”），不是 $\dd s^2$ —— 因为它是 T-对偶之后（$\tphi$ 一侧）的度规；后文 (5.14)、(5.20) 沿用同一记号。]
